% --- OBJETIVOS ---
\section*{Objetivos}
\begin{itemize}[label=\textbullet]
    \item \textbf{General:} Comprender y aplicar las leyes fundamentales de la termodinámica para analizar el comportamiento de los sistemas energéticos y la transformación de la energía.
    \item \textbf{Específicos:}
    \begin{itemize}
        \item Diferenciar y analizar los mecanismos de transmisión de calor (conducción, convección y radiación).
        \item Modelizar el comportamiento de los gases perfectos y comprender la implicancia de la Ley de Avogadro.
        \item Aplicar el Primer Principio de la Termodinámica en diversos procesos y transformaciones.
        \item Analizar la eficiencia de los ciclos termodinámicos, comprendiendo los límites teóricos impuestos por el Ciclo de Carnot.
        \item Estudiar los procesos de cambio de fase, específicamente la vaporización, y su aplicación industrial.
    \end{itemize}
\end{itemize}

% --- METODOLOGÍA ---
\section*{Metodología}
\begin{itemize}[label=\textbullet]
    \item \textbf{Enfoque:} Teórico-Práctico. Se priorizará la resolución de problemas reales y la aplicación de modelos matemáticos a fenómenos físicos.
    \item \textbf{Estrategias:}
    \begin{itemize}
        \item Clases expositivas con soporte visual para la conceptualización.
        \item Resolución de guías de trabajos prácticos.
        \item Análisis de casos y simulaciones de ciclos termodinámicos.
        \item Fomento del pensamiento crítico mediante la deducción de fórmulas a partir de principios básicos.
    \end{itemize}
\end{itemize}

% --- EVALUACIÓN ---
\section*{Evaluación}
\begin{itemize}[label=\textbullet]
    \item \textbf{Criterios:}
    \begin{itemize}
        \item Capacidad de análisis y síntesis de los conceptos teóricos.
        \item Precisión en la aplicación de fórmulas y resolución de problemas numéricos.
        \item Cumplimiento y entrega en tiempo y forma de los trabajos prácticos.
        \item Participación activa y capacidad de argumentación técnica en clase.
    \end{itemize}
    \item \textbf{Instrumentos:}
    \begin{itemize}
        \item Evaluaciones escritas parciales.
        \item Entregas de trabajos prácticos corregidos.
        \item Instancias de defensa oral de los ejercicios resueltos.
    \end{itemize}
\end{itemize}